\documentclass[11pt,a4paper]{article}
\usepackage{amsmath,amssymb,amsthm}
\usepackage[margin=2.5cm]{geometry}
\usepackage{hyperref}

\newtheorem{definition}{Definition}[section]
\newtheorem{theorem}[definition]{Theorem}
\newtheorem{proposition}[definition]{Proposition}
\newtheorem{remark}[definition]{Remark}
\newtheorem{conjecture}[definition]{Conjecture}

\title{Hyperseed Formalization:\\
  Paraconsistent Interzones}
\author{ProtomegaTron (automated formalization)}
\date{2026-09-18}

\begin{document}
\maketitle

\begin{abstract}
We formalize Ben Goertzel's ``Paraconsistent Interzones'' (2021-08-13)
within the Hyperseed ontology. Key mappings: interzones as fiber overlap
regions where contradictions flourish, creativity as interzone navigation,
AGI as native interzone dweller, and control systems as consistent fiber
that breaks down in interzones.
\end{abstract}

\tableofcontents

\section{Fiber overlap: the interzone}

\begin{theorem}[Interzone = fiber overlap --- claims 1--4, 18, 21]
An interzone is a region of the base space where fibers from different
bundles overlap and contradict:
\[
  \text{Interzone}(b) = \{b \in B : F_1(b) \wedge F_2(b)
  \wedge F_1(b) \text{ contradicts } F_2(b)\}
\]
The interzone belongs to neither territory --- it has its own emergent
topology, distinct from both contributing territories. Contradictions
in the interzone are resources, not errors: they provide the raw material
for novel synthesis.

Interzones are liminal spaces --- thresholds where established fiber
structures break down and new ones can emerge. The most creative and
intellectually productive work happens in these spaces.
\end{theorem}

\section{Classical failure, paraconsistent success}

\begin{proposition}[Paraconsistency enables interzones --- claims 8--11, 22]
Classical logic fails in interzones because contradictions cause explosion:
\[
  F_{\text{classical}} \models (A \wedge \neg A) \implies
  F_{\text{classical}} \models \psi \;\forall\psi
\]
Paraconsistent fiber can operate where classical fiber cannot:
\[
  F_{\text{paracons}} \models (A \wedge \neg A) \quad \text{but} \quad
  F_{\text{paracons}} \not\models \psi \;\forall\psi
\]
Control systems require consistent fiber and break down in interzones.
Paraconsistent fiber can grade contradictions and reason productively
about them, extracting useful conclusions from contradictory premises.
\end{proposition}

\section{Interzone navigation: creativity}

\begin{theorem}[Creativity = interzone navigation --- claims 2, 12--14, 19]
Creative cognition is navigation through interzones --- finding novel
fiber structure where existing fibers clash:
\[
  F_{\text{new}}(b) = \text{navigate}(F_1(b), F_2(b))
  \quad \text{where } F_1(b) \text{ contradicts } F_2(b)
\]
The new fiber resolves (but does not eliminate) the contradiction ---
it provides a perspective that integrates insights from both
contradictory perspectives while transcending both. This parallels
interdisciplinary scientific discovery, mathematical creativity at
subfield boundaries, and cultural synthesis at cultural borders.
\end{theorem}

\section{AGI as interzone dweller}

\begin{conjecture}[Native paraconsistency --- claims 15--17, 20]
AGI should be natively paraconsistent --- dwelling comfortably in
interzones rather than retreating to safe consistent subspaces:
\[
  \text{AGI} \in \text{Interzone}(F_1, F_2, \ldots, F_n)
  \quad \text{natively, not as exception}
\]
Multi-framework reasoning means maintaining multiple contradictory
perspectives simultaneously and synthesizing them --- not choosing
one framework and discarding others. The AGI's home is the interzone.
\end{conjecture}

\section{Cross-references}

\begin{remark}[Connections]
\begin{itemize}
\item \textbf{The Mind-Blowing Paraconsistency} (2021-09-23): deep dive
  into the paraconsistent logic that enables interzone reasoning.
\item \textbf{Symbolic Ruminations} (2022-08-06): the symbolic/non-symbolic
  boundary as an interzone.
\item \textbf{Parafinity} (2022-01-15): the finite/infinite boundary as
  a parafinite interzone.
\item \textbf{Open-Ended Motivations} (2021-08-13): interzone navigation
  as a form of open-ended growth.
\end{itemize}
\end{remark}

\end{document}
